\clearpage
\begin{summarybox}
  \textbf{\textcolor{CSummary}{一句话记住}}

  \textbf{把反复出现的根号当作新字母，把陌生的式子改成熟悉的式子；但原题允许取到的数，一个也不能多、一个也不能少。}

  \medskip
  \textbf{\textcolor{CSummary}{四类常见根式，优先这样想}}
  \begin{itemize}[leftmargin=*,itemsep=0.45em]
    \item $\sqrt{mx+n}$（$m\ne0$）：设 $t=\sqrt{mx+n}\ge0$，得到 $x=\frac{t^2-n}{m}$。
    \item $\sqrt{a^2-x^2}$（$a>0$）：设 $x=a\sin\theta$，角度限制为 $[-\frac\pi2,\frac\pi2]$，根号变成 $a\cos\theta$。
    \item $\sqrt{x^2+a^2}$（$a>0$）：可用 $x=a\tan\theta$；若式子里同时有 $x$ 与根号的加减，先试平方差公式。
    \item $\sqrt{x^2-a^2}$（$a>0$）：先分 $x\ge a$、$x\le-a$；若出现 $x-\sqrt{x^2-a^2}$，先乘 $x+\sqrt{x^2-a^2}$。
  \end{itemize}

  \medskip
  \textbf{\textcolor{CSummary}{最值得记的两条快算式}}
  \[
    \sqrt{x^2+a^2}-x=\frac{a^2}{\sqrt{x^2+a^2}+x}\quad(a>0),
  \]
  \[
    x-\sqrt{x^2-a^2}=\frac{a^2}{x+\sqrt{x^2-a^2}}
    \quad(x\ge a>0).
  \]

  \medskip
  \textbf{\textcolor{CSummary}{考场上的四步检查}}
  \begin{enumerate}[leftmargin=*,itemsep=0.35em]
    \item 原来的 $x$ 可以取哪些值？
    \item 换成 $t$ 或 $\theta$ 后，范围究竟是什么？
    \item 新式子怎样最省事地求最大值、最小值或上下界？
    \item 上下界能取到吗？中间的值能全部取到吗？
  \end{enumerate}

  \textbf{最后提醒：} \textbf{$\sqrt{u^2}=|u|$；范围先行，化简随后。}
\end{summarybox}
